\begin{exercise}[Allotropie du fer]
L'évolution des paramètres de maille a été obtenue à différentes températures :

\begin{center}
    \begin{minipage}{0.6\textwidth}
    \centering
    \begin{tabular}{cccccc}
    \toprule
    Température & [\si{\celsius}] & 20 & 912 & 1394 & 1535 \\ 
    \midrule
    \ce{Fe(\alpha  )} & [\si{\angstrom}] & 2.86 & 2.89 & - & - \\ 
    \ce{Fe(\gamma  )} & [\si{\angstrom}] & - & 3.63 & 3.67 & - \\ 
    \ce{Fe(\delta  )} & [\si{\angstrom}] & - & - & 2.90 & 2.91 \\ 
    \bottomrule
    \end{tabular}
    \end{minipage}\begin{minipage}{0.4\textwidth}
    \captionof{table}{Mesure expérimentale de quelques paramètres de maille $a$ du fer à différentes températures et sous différents allotropies.}
    \end{minipage}
\end{center}

\begin{questions}
    \question Proposer une méthode expérimentale permettant d'acquérir ces paramètres.

    \question En supposant que pour un allotrope donné, le paramètre de maille évolue de manière linéaire avec la température, représenter son évolution avec la température.

    \question En déduire l'évolution du volume avec la température.

    La pression a une influence sur la stabilité, en température, des allotropes/phases. Lorsque la pression augmente, on constate que : 
    \begin{align*}
    P \nearrow \implies T_{\alpha \leftrightarrow \gamma}^\star (P) \searrow \\
    P \nearrow \implies T_{\gamma \leftrightarrow \delta}^\star (P) \nearrow \\
    \end{align*}

    \question Justifier qualitativement ce comportement à l'aide de la question précédente.
\end{questions}
\end{exercise}